\chapter{Introduction and Principles}

\section{Introduction}

The initial concept of this book was to write a list of species with a one line
description of the best method for germinating the seed. Such a book would be
patterned after the \textit{The Seedlist Handbook} of Bernard Harkness (ref. 1) and
would supplement that book. After reading the six latest books on seed
germination (refs. 2-7), it was evident that the book could not be constructed
from data in the literature. The decision was made to embark on a broad program
of studying seed germination. At this time the realization came that the
germination of seeds is a chemical process, and seed germination should be
studied with the same techniques and logic used to study chemical processes.
Perhaps these techniques, which had so drastically changed our concepts of
chemical processes, would have a similar effect on our concepts of seed
germination.

Initially the studies were restricted to alpine and rock garden plants, but with
each new discovery, the thrill of the hunt increased. Ultimately the studies were
broadened to include trees and shrubs and finally tropical plants and the grasses. It
became evident that the work was of importance not only to plant growers but also to
research, biologists and research chemists. The question arose whether to write
separate books for each of these three audiences or to attempt a tour de force putting
it all in one book. The tour de force was chosen, and this has necessitated some
words of explanation directed at each of the three audiences.

For the plant grower, the directions for optimum germination of nearly 2500
species will be of the most interest. These directions are summarized in Chapter 20
and arranged in alphabetical order of the genera. A number of abbreviations are used
in Chapter 20, and these are all explained at the beginning of that chapter. \textit{These
directions for optimum germination can be followed without reference to an y other cart
of the book.} After the seed has been germinated, there is the problem of how to best
handle the small seedlings. This subject is taken up in Chapter 14, and several highly
efficient methods are described.

I can appreciate that the first reaction to the detailed recommendations and
extensive data is that it is all too complicated. There will be the temptation to go back
to old ways, to water long rows of pots twice a week for months on end, to expect only
a small percentage to germinate, and to resort to growing the easy ones dismissing
the rest as difficult and recalcitrant.

There are better ways with which you can expect success rates of 90\% and
better. Best of all it can be done with a \textit{fraction of the effort and expense
and most of all a fraction of the space} that are used in traditional methods.
The procedures require the usual pots and media plus high wet strength paper
towels and polyethylene sandwich bags called Baggies (both available at your
grocery store) and permanent ink pens (available at any office supply store).

The procedures are described in detail in Chapters 3 and 14. As you become more
comfortable with these new techniques, you will have the confidence to try
anything. Especially become familiar with germinating seeds in pots enclosed in
polyethylene bags. Plants can be grown by this procedure for six months to a
year under fluorescent lights with watering reduced to less than once every two
months. Also become familiar with germinating seeds in moist paper towels
enclosed in baggies. For seeds that require extended treatments, these lengthy
periods of time can be spent in the moist paper towels. Watering is reduced to
once every several months, the towels occupy very llittle space, and no light
is needed. Be sure to read Chapters 3 and 14 to understand that the
polyethylene bags must not be sealed too tightly and that thin polyethylene
film is better than thick films. Aerobic conditions must be maintained.

Some species require gibberelic acid for germination, but do not let that frighten
you. It is a natural requirement, and a highly efficient procedure has been developed
for applying the gibberelic acid (see Chapters 3 and 12). The only additional materials
that are needed are toothpicks and the gibberelic acid. No weighings or volume
measurements are needed. The gibberelic acid is used so efficiently that one gram is
sufficient for a thousand experiments since only one cubic millimeter is required for
each experiment.

So now you have the program. Sit down on a cold winter's night with the list of
seeds that have been ordered or perhaps have already arrived. Arrange them in
alphabetical order. Then go through Chapter 20 looking up each one. Arrange them
in groups that require similar treatments. Plan your strategies. Time the operations so
that the seedlings started indoors are ready to be set outdoors at appropriate times.
Use the highly efficient short cuts described in Chapter 3 and 14.

Let me encourage the plant grower to read the other chapters. \textit{The rate studies
are framed in language that everyone can understand.} A rate plot is simply a graph of
the number or percent seeds germinated on the vertical axis against time on the
horizontal axis. Such a plot with its induction time and rate curve contains much more
information than percent germination and is a more sensitive way to determine the
effect of variables. \textit{Everyone can understand the first principal of mechanistic
chemistry.} If a treatment A or conditions A are required to get germination under
conditions B, critical proceses were taking place during A and these were likely taking
place at their optimum rates (greatest speeds). Thus if a period of three months at 40
deg F is required to get germination at 70, the period at 40 was a time when something
was happening (typically destruction of germination inhibitors); and it is misleading to
refer to this period as a period of dormancy or a period of breaking dormancy.

For \textbf{biologists} the rate studies and rate theory in Chapters 4 and 19 should
encourage all biologists to make more use of rate studies. One of the principles of
physical organic chemistry is that complex multiple step processes can obey simple
rate laws. Generally in a complex sequence, one step will be the slow or rate
determining step. This step not only dictates the overall rate; but it also dictates the
rate law for the overall process. Chapter 19 shows the excellent fit of some
germination rates to simple rate laws.

Biologists will also be interested in the special ability of biological-systems
to show enormous changes in rate with temperature. These can be (a) the result
of configurational changes in proteins from enzymatically active configurations
to inactive ones or (b) the result of structural changes in membranes that
alter their permeability.  These changes can be complete over a temperature
range of as little as five degrees.  Many examples of enormous rate changes are
described in Chapters 6 and 7, and the theory is discussed in Chapter 19. I
find that biologists are quite familiar with the extreme sensitivity of
metabolism to temperature changes, but they do not always look on these as
rates and rates with enormous temperature coefficients.

For the \textbf{chemist} the ability of biological systems to show enormous changes in
rate with temperature opens up new rate behaviors that are unknown in small
molecule chemistry. Examples are the germinations that take place under conditions
of oscillating temperatures. The theory of these is discussed in Chapter 19. Do not be
put off by the simplicity of the experiments. At this stage the basic phenomena have to
be discovered, and the simple experiments are best suited for this purpose. In my
career we used every form of spectroscopy, and these were available. It will be for
future workers to Use these more sophisticated methods for determining the intricate
chemistry involved.

This book depends on botanical taxonomy and the concept of a species.  Linneaus
believed that species were immutable. Darwin and the principle of evolution (it
is not a theory anymore) changed all that. Fortunately there is a solution to
this conflict. The solution has been presented in its clearest form by Stephen
Gould in his article \textit{What is a Species}. His views are adopted and are described
in Chapter 16.  The experiments, the writing, the publication, and the
distribution of this book were done by myself. The book has been distributed
\textit{below cost}. The effort was pure pleasure as it allowed me to combine my career
in chemistry where I had published one hundred and fifty papers with a
secondary career in in horticulture where I had published twenty papers.

This book and the experiments described would not have been possible without
the contributions of seeds from over one hundred people, many botanic gardens, and
the seed exchanges of horticultural societies. These are all acknowledged in Chapter
23. Since writing the first edition, many suggestions have been received (and some
sharp criticisms). These have all been of value and all had some influence on this
second edtion. The people who sent these are acknowledged in Chapter 23. Also
acknowledged are the many wonderful teachers that I had at all levels and the many
wonderful colleagues and graduate students with whom I associated.

\section{Principles}

\begin{itemize}
\item \textsc{Every species of plant has one or more mechanisms
for delaying germination until the seed is dispersed}
(Chapters 5-12). This principle pervades every aspect of seed germination. Without
such mechanisms, seeds would germinate, in the capsule or fruit and would never be
dispersed. The challenge in germinating seeds is to overcome these delay
mechanisms. It is proposed to call all treatments used to overcome delay mechanisms
as \textit{conditioning}. With 95\% of the species studied, the delay mechanism is chemical.
With the remaining 5\% it is physical. Although some species have seeds that
germinate qiuckly after dispersal such as Oronticum, Populus, and Tussilago; they still
must have a mechanism to block germination before dispersal. A few species (the
mangrove and Polygonum viviparum) germinate before dispersal, but the germination
is so limited that in effect ungerminated seeds are being dispersed.

\item \textsc{Temperatures of 40 or 70 deg. F were used for both germinations and
conditioning and rarely was any other temperature needed} (Chapter 3). A few
species germinated better at temperatures in the range 90-100 or after brief
exposures to even higher temperatures (Gomphrena, Kalmia, and Mimosa in Chapter
20). Such behaviors were rare. Although it is intended to conduct exploratory
experiments on temperatures between 40 and 70, it is not expected that these
will alter in any major way either the principles or the recommended
procedures.

\item \textsc{Members of the same family, the same genus, and even closely related
species may have different mechanisms for delaying germination.} The data in
Chapter 20 are arranged by genera, but this is not meant to imply any
generalized behavior within each genus. While there are trends, behaviors are
as much a result of adaptations to the environment in which the plant grows as
by evolutionary associations.

\item \textsc{Fungal products of the gibberelin type have a major effect on germination
of at least 25\% of all species} (Chapters 3 and 12). Some species have evolved gibberelins as a requisite for
germination. This has important survival value, particularly for small plants in cold
desert areas. With other species gibberelins may have a large effect on germination,
but germination can be effected without it. It is remarkable that the species that
required light for germination usually germinated in the dark with gibberelin. The
reasons for this relation are not understood. Gibberelins can also be without effect or
deleterious to germination so that gibberelins should not be conceived as some "cure
all" for recalcitrant germination. The effects are complex and an example was found
where gibberelic acid-3 affected the second step but not the first step in a two step
germinator (Lilium canadense). Rarely germination with gibberelic acid-3 takes place
at 40 (the rosulate violas) instead of at 70.

\item \textsc{About 50\% of temperate zone plants use a delay
mechanism that is destroyed by drying} (Chapter 5). The prevalence
of this mechanism leads to a common misconception that all seeds will germinate
when exposed to warmth and moisture. The commmon practice of regarding such
seeds as "immediate germinators" obscures the fact that a drying period is needed
before the seed will germinate. The length of time needed for the drying is dependent
on the species, the temperature; and the relative humidity. During the drying chemical
changes take place that are amenable to rate analysis.

\item \textsc{Most species have at least two delay mechanisms,
one being in the nature of a chemical time clock} (Chapter 4).
Most seeds have an induction period of one to eleven weeks after moistening and
before germination begins. This is true even after any requisite drying or any requisite
cold or warm moist cycles. In the past such induction periods have been regarded as
time required for physical transfer of oxygen and water or time required for growth of
the embryo. However, such explanations would generate a rate plot where there was
a \textit{gradual} onset of germination. What is generally found is that there is a \textbf{sudden} onset
of germination which is more typical of a chemical time clock.

\item \textsc{It is common for the delay mechanism to be
destroyed at one temperature followed by germination at
another temperature} (Chapter. 6). It has long been recognized that many
seeds require a period around 40 deg. F followed by a shift to 70 deg. F before
germination will occur. In the present work it was found that there were nearly as many
examples of the reverse, namely a period of time at 70 followed by germination at 40.
Both behaviors require that the processes involved in conditioning have enormous
changes of rate with temperature.

\item \textsc{Some species have several delay mechanisms that
must be destroyed under different conditions of
temperature and time and in sequence} (Chapter 7).

\item \textsc{Some species germinate under conditions of
oscillating temperatures} (Chapter 10). Such seeds are programmed to
germinate in the fall or spring when daily temperatures oscillate widely. Such
oscillating temperatures were simulated in the in the laboratory with Cleome, and the
same stimulation of germination resulted. The frequency of the temperature
oscillations were varied with little effect on the rate plots. iThis was interpreted as
showing that the results were due.to the oscillating temperatures and not germination
at some temperature intermediate between 40 and 70. Experiments are planned to
investigate these intermediate temperatures and settle the question decisively.

\item \textsc{Light is an important variable and is a requirement
for germination in some species} (Chapter 11). The requirement of light
for germination is characteristic of plants growing in swamp's or woodland where the
availability of light is more of a problem than the availabilty of water. The requirement
of light for germination can usually (but not always) be replaced by treatment of the
seeds with gibberelic acid-3 as mentioned in number 4 above. Rarely light can block
germination, and the survival value of this is discussed. A photosynthesizing seed
coat was a factor in one species (Hymenocallis) and a photosynthesizing seed
capsule was a factor in another (Zantedeschia).

\item \textsc{Seeds in fruits often have chemical inhibitors in
the flesh of the fruit that block germination and must be
removed by washing before germination will take place}
(Chapter 8). Usually these chemical inhibitors are water soluble so that soaking in
water with daily rinsing for seven days is sufficient for their removal. Less commonly
the inhibitors are oil soluble so that the washings must be conducted with aqueous
detergents. The inhibitors are continually diffusing into the seed, and the seed is
continually destroying the inhibitor. The rates of the destruction of inhibitors inside the
seeds can be determined. Sometimes this rate is so fast that the seeds germinate in
the third day of washing. Despite the physically hard nature of many seeds of this
type, there is a cleft in the seed coat so that producing a hole in the seed coat has no
favorable effect. (Prinsepia is an exception). Although the above mechanism if often
found in seeds in fruits, other mechanisms involving light and dry storage
requirements are also found.

\item \textsc{Some species have two step germinations} (Chapter 6).
In the first step the radicle emerges, and a small bulb or corm with root is formed. In
the second step a photosynthesizing leaf or leaves are formed. Usually the first step
takes place at 70. Then a period of time at 40 is required followed by a return to 70.
Sometimes simply an extended period of time at 70 or 40 serves to initiate the second
step. The concept of two step germination should not be confused with epigeal and
hypogeal germination. Epigeal (photosynthetic cotyledons) and hypogeal (aborted
cotyledons) can be either one step or two step germinators.

\item \textsc{Seeds of species from cold desert areas
generally germinate at low temperatures, typically 35-40
deg. F} (Chapters 6 and 7). This is a principle of general validity. The survival value
is obvious. Seeds of plants from cold desert areas have to get a quick start in spring
when moisture is still available.

\item \textsc{Rate data are less ambiguous than the traditional
percent germination and are more sensitive to internal
changes in the seeds.} The traditional percent germination can mean
anything from the percent viable seeds that germinate to the percent of normal size
seed coats that contain viable seeds. It will depend on the degree of sophistication
used in detecting viable seeds from both non-viable seeds and empty seed coats.

\item \textsc{Some species produce quantities of normal size seed coats which are empty}
(Chapter 3). This has survival value. The calculation of percent germinations
is complicated by this factor, but rates of germination are not.

\item \textsc{Results that appear to be inconsistent can arise
from several causes} (Chapter 3). Causes include the presence or absence
of gibberelin producing fungi in the media, variations in storage of the seed,
uncertainties in the identification of the botanical species and subspecies, and
possibly variations in the climate in different years or in different geographical
localities. Pathogenic fungi were not a problem in the procedures used, and the
evidence for this is presented in Chapter 3.

\item \textsc{Parasitic plants may require exogenous
chemicals from the host for germination, but usually they
do not} (Chapter 12). Although exogenous chemicals may be required for the later
development, the germination of the seed usually does not have such requirements.

\item \textsc{Germination of seeds usually follow zero order
or first order kinetics with good precision} (Chapters 4 and 19).
The reasons why complicated biological processes can follow simple mathematical
equations is explained in Chapter 19.

\item \textsc{Seeds of a significant number of species are sold
and distributed by commercial seed companies in a dod
(dead on delivery) state} (Chapter 13). Many species have seeds that do
not survive very long in dry storage even at 40. Ways to store and distribute seeds in
moist states are described in Chapter 13.

\item \textsc{The concept of dormancy and requisite cold
cycles has been completely revised} (Chapter 19).
\end{itemize}

